% Full proofs for the explicitly selected binary, unit-free net calculus.
\section{Multiplicative structures, splitting, and cut}

\begin{annexitem}{annex:p07.mll.desequentialization}
 {thm:p07-mll-desequentialization}{All construction and component-count cases}
\begin{proof}
A severed premise arc retains a new boundary endpoint, so it stays in its
original component. An axiom is a two-leaf star and a hypothesis is one
loose edge; both are trees. Nullary Mix is empty. Exchange changes nothing.
For tensor the premise structures are disjoint. Joining their selected
conclusions through the new tensor joins different components, creates
no cycle, and decreases the sum of component counts by one. Two trees
give a tree. At par the chosen premise is extended through the new
vertex and conclusion; the other premise is left loose. This is a
pendant extension in the selected switching, creating no cycle and
changing no component count. Binary Mix is disjoint union.
Induction gives both the acyclic and Mix-free tree assertions.

The component formula holds in every case. Axiom and hypothesis give
\(1\); nullary Mix gives \(0=1-1\); par leaves it unchanged.
Tensor adds the two premise formulas and subtracts one, whereas binary
Mix adds them and accounts for its additional \(\mathsf{Mix}_2\).
Thus \(d=1+\#\mathsf{Mix}_2-\#\mathsf{Mix}_0\), independently of switching.

Induct on the proof receiving a substituted hypothesis. At that leaf
graph gluing is exactly replacement. At par it commutes with adding
the vertex; at tensor or Mix it occurs in just one disjoint premise.
The carried conclusions remain loose. Hence identifying the output
arc with the marked input arc gives exactly the substituted proof's graph.
\end{proof}
\end{annexitem}

\begin{annexitem}{annex:p07.mll.trip-translation}
 {thm:p07-mll-trip-translation}{Every original rule preserves a single orbit}
\begin{proof}
The axiom, terminal, tensor and par rows partition exit states by their
unique producer and their consumer or terminal status. In those rows
every state has exactly one successor and predecessor; the successor
is a permutation of the \(2p\) states.

An axiom has the one orbit
\(A^+,(A^\perp)^-,(A^\perp)^+,A^-\), then returns to \(A^+\).
Exchange changes no incidence. At a par with new root \(C\), setting
\(L\) replaces the old terminal return \(A^-\mapsto A^+\) by
\(A^-\mapsto C^-\mapsto C^+\mapsto A^+\), leaving the other
terminal return unchanged. The one old orbit acquires the two new states.
Setting \(R\) does the same at \(B\).
At tensor the two premise orbits are disjoint. With \(C=A\otimes B\),
setting \(L\) replaces the terminal returns by
\[
 A^-\mapsto C^-\mapsto C^+\mapsto B^+,\qquad B^-\mapsto A^+.
\]
This joins the two old orbits and inserts both new states. Setting \(R\)
interchanges the premises. The induction applies to arbitrary settings
at the old links, so every setting has one long trip.
\end{proof}
\end{annexitem}

\begin{annexitem}{annex:p07.mll.paths}
 {thm:p07-mll-paths}{Local/global and cycle-extraction facts}
\begin{proof}
An internal vertex of a simple cycle occurs only once. Visiting both
premises of a par there makes a bridge. The only repeated vertex can
be the endpoint; a strong cycle based at that par starts through its
conclusion and cannot use both premises at the endpoint. Thus a strong
bridge-free simple cycle uses at most one premise of every par globally.
Complete those choices at unvisited pars to obtain a switching containing it.

A bridge at a concatenation junction could occur only at a par.
The second strong path starts through its conclusion, not a premise,
so the junction is not a bridge.
A simple strong-weak path has no immediate arc reversal. At a junction
its final arc is a par premise and the next path's first arc is that
par's conclusion; those arcs differ. Thus the concatenation has no
immediate reversal. In a nonempty closed concatenation choose a shortest
nonempty closed contiguous subwalk. It has no repeated internal vertex,
since otherwise a shorter closed subwalk could be selected. A repeated
arc would repeat an internal endpoint, except for a two-step immediate
reversal, which has been excluded. It is a simple cycle.
Parallel arcs remain distinct throughout this argument.
\end{proof}
\end{annexitem}

\begin{annexitem}{annex:p07.mll.bungee}
 {thm:p07-mll-bungee}{The minimal-bridge cycle cannot admit the bypass}
\begin{proof}
Let \(\kappa\) be the bridge pier. Truncate the bypass \(\gamma\) at
its first new encounter with \(\omega\), so the paths meet only at
\(\kappa\) and its target \(u\). If \(u=\kappa\), \(\gamma\) is
already a strong bridge-free simple cycle and hence a switching cycle.

Otherwise orient \(\omega\) so \(u\) is on the portion from \(\kappa\)
to its ending occurrence of \(v\). If \(u=v\), retain the starting
portion that begins with a conclusion of \(v\). Replace the portion
from \(\kappa\) to \(u\) by \(\gamma\). The resulting closed path
\(c\) is simple, since the bypass has no other vertex on \(\omega\).
It contains a conclusion of \(v\): for \(u\ne v\) the retained portions
contain the two arcs at the based occurrence of \(v\), one of which is
the specified conclusion; for \(u=v\) the chosen starting arc supplies it.

The old bridge at \(\kappa\) disappears, since \(\gamma\) starts there
through a conclusion. No bridge occurs inside \(\gamma\).
Only \(u\) can add a bridge. Minimality of the number of bridges in
\(\omega\) now forces a new bridge at \(u\), and forces the removed
portion from \(\kappa\) to \(u\) to have no other bridge.

Follow \(\gamma\) and then that removed portion backwards.
This is a simple cycle, strong at \(\kappa\), with bridge-free interiors.
At \(u\), the new bridge in \(c\) uses the arriving bypass premise and
the retained outgoing premise. The other original arc at \(u\), used
by the backwards portion, must therefore be the par's conclusion.
Hence no bridge occurs at this junction. At \(\kappa\) the bypass
departure is a conclusion, so again there is no bridge.
The resulting strong bridge-free cycle is a switching cycle.
\end{proof}
\end{annexitem}

\begin{annexitem}{annex:p07.mll.splitting}
 {thm:p07-mll-splitting}{The finite dependency argument produces the splitter}
\begin{proof}
The relation \(\prec\) is irreflexive because a bridge pier is internal
in a simple cycle based at its source. Suppose a cyclic dependency
\(v_1\prec v_2\prec\cdots\prec v_n\prec v_1\) exists.
For each edge choose its minimal-bridge cycle \(\omega_i\), starting
with \(v_i\)'s conclusion and having first bridge at \(v_{i+1}\).
Its prefix \(\gamma_i\) ending there is simple, strong-weak and bridge-free.
The closed concatenation is bridge-free and contains a simple cycle
\(\sigma\) as a contiguous subpath by the path lemma.

That cycle cannot be contained in one prefix, which is simple and open.
Start the cyclic listing at the prefix containing the start of \(\sigma\).
Let \(v_{i+1}\) be the first junction strictly inside \(\sigma\).
Its suffix to the end of \(\sigma\) starts along the next prefix's
conclusion, is simple and bridge-free, and ends on \(\gamma_i\),
hence on \(\omega_i\). Since \(v_{i+1}\) is a bridge pier of
\(\omega_i\), the bypass lemma gives a switching cycle.
Acyclicity excludes the cyclic dependency. The transitive closure
of \(\prec\) is consequently a strict partial order.

Choose a maximal member \(v\) of the finite nonempty set of pars.
If \(v\) were not splitting, choose a cycle containing its conclusion
with the minimum number of bridges, starting with that conclusion.
If there is no bridge it is a strong bridge-free switching cycle.
Otherwise its first bridge pier \(w\) gives \(v\prec w\).
Both alternatives contradict respectively acyclicity or maximality.
Thus \(v\) is splitting. No sequentialization or tensor splitting
theorem has been assumed.
\end{proof}
\end{annexitem}

\begin{annexitem}{annex:p07.mll.sequentialization}
 {thm:p07-mll-sequentialization}{Internal-par factorization and all base cases}
\begin{proof}
Induct on the sum of vertices and arcs, allowing marked hypotheses.
If a par exists, choose the proved splitting par with premises \(A,B\)
and conclusion \(C=A\bindnasrepma B\). Delete its vertex but keep
the three arcs with their former incidences undefined.
The conclusion side is separate from the premise sides: an outside
connection to either premise side would complete a simple cycle
through its conclusion. Let \(N_1\) contain the premise sides and
\(N_2\) the conclusion side. Allocate unrelated components once only
to either. Then \(N_1\) has conclusions \(A,B\), while \(N_2\)
has the distinguished marked hypothesis \(C\).
Each is smaller and has no new switching cycle.
Sequentialize them by induction, apply par to the proof of \(N_1\),
and substitute it for the distinguished hypothesis of the proof of
\(N_2\). The proved substitution property restores every original arc.

With no pars, every cycle would be a switching cycle, so the graph is
a forest. If there is a tensor, follow its conclusion through consuming
tensors. Formula complexity increases and the structure is finite,
so a terminal tensor is reached. Delete it and its terminal output.
Its two premise sides are in different components, otherwise there
was a cycle through that tensor. Allocate other components to either
side once only, sequentialize the two smaller structures, and apply
tensor to restore it.

With no pars or tensors, every vertex is a terminal axiom.
Delete one axiom and its outputs, sequentialize the remainder, and
Mix with the axiom proof. With no vertices but an arc, it has both
endpoints undefined; remove it, sequentialize the remainder and Mix
with that marked hypothesis. The empty case is nullary Mix.
Every recursive call decreases the finite measure, and the construction
retains exactly the original marked hypotheses and boundary occurrences.
\end{proof}
\end{annexitem}

\begin{annexitem}{annex:p07.mll.connected-sequentialization}
 {thm:p07-mll-connected-sequentialization}{Eliminating the auxiliary Mix rules}
\begin{proof}
The forward construction direction was proved already. Conversely,
sequentialize the closed connected structure; closedness gives no
marked hypothesis leaves.

An empty-sequent proof contains only Mix: axiom, tensor and par have
nonempty conclusions; binary Mix has empty conclusion only if both
premises do. Recursively collapse any empty proof to one nullary Mix.
In a proof of a nonempty sequent, discard a binary Mix with an empty
premise, retaining the other. Recurse into all other premise proofs.
A nullary Mix cannot supply a tensor's required principal formula or
a par's two principal formulas. This finite recursive pruning preserves
the graph and leaves a nonempty proof without nullary Mix.
Its component count is \(1+\#\mathsf{Mix}_2\).
Since it is one, no binary Mix remains. The proof is closed and
Mix-free, establishing the exact converse.
\end{proof}
\end{annexitem}

\begin{annexitem}{annex:p07.mll.cut-preservation}
 {thm:p07-mll-cut-preservation}{The two-switch outside-forest calculation}
\begin{proof}
Fix all par settings outside a chosen redex.
For an axiom contraction its axiom and cut form a path between the
opposite premise producer and the companion output's consumer or
boundary. If an outside par severed that companion incidence, use
its new boundary endpoint. The contraction shortens that path with
the same endpoints. The same axiom cannot produce both cut premises:
the two incidences would be a switching cycle. Hence this is legitimate
path shortening, preserving acyclicity and components in every setting.

For a principal cut, remove its tensor, par and cut and the two internal
conclusion arcs, freeing the four premise ports. The remaining outside
graph is a forest \(Q\). Call its ports \(a,b,a',b'\).
The old left setting joins \(a,b,a'\) by a tree; those ports must be
in three distinct components of \(Q\), or it creates a cycle.
The old right setting likewise requires \(a,b,b'\) in three distinct
components. Thus the only pair of the four ports that can already be
in the same component is \(a',b'\).
The reduct joins \(a\) to \(a'\) by one cut and \(b\) to \(b'\) by
another. Each is a join between different components, including after
the first join, and therefore neither creates a cycle.

Either old setting reduces the component count of \(Q\) by two.
The reduct's two joins do the same. Other components stay unchanged.
Thus every outside setting has the same number of components after
contraction, and connectedness is preserved when that number is one.
No sequentialization or unproved splitting dependency is used here.
\end{proof}
\end{annexitem}

\begin{annexitem}{annex:p07.mll.cut-termination}
 {thm:p07-mll-cut-termination}{The strictly decreasing integer and all cut producers}
\begin{proof}
A principal contraction removes three vertices and inserts two, reducing
their count by one. An axiom contraction removes the axiom and cut,
reducing it by two. Hence there is no infinite reduction.

In a closed structure every cut premise has an axiom, tensor or par
producer. An axiom producer admits the axiom contraction; correctness
excludes the exceptional same-axiom loop. Otherwise dual outer
connective labels force a tensor opposite a par, admitting the
principal contraction. Thus every irreducible correct closed structure
is cut-free. Preservation keeps these hypotheses valid at every step.

A nonempty closed correct structure has positive switching component
count, preserved by contraction. Its final reduct cannot be empty.
A closed vertex-free structure has no arc, since every arc needs a
producer; hence the final reduct retains at least one vertex.
From \(n>0\) vertices there are at most \(n-1\) contractions.
The empty structure takes zero.
\end{proof}
\end{annexitem}

\begin{annexitem}{annex:p07.mll.cut-confluence}
 {thm:p07-mll-cut-confluence}{The overlap cases and well-founded confluence}
\begin{proof}
Two different cut redexes have disjoint internal vertices except possibly
an axiom. A tensor or par has one output with at most one consumer,
so cannot belong to two cut redexes. Disjoint surgeries commute:
each retargets only its specified incidences, retaining the producer
and consumer required by the other.

If two cuts share an axiom, they use its two different outputs.
Deleting either cut and that axiom leaves one cut joining the two
former opposite premises. The two outcomes have the same formula
labels and outside endpoints and are isomorphic, with the cut's
symmetric order and internal vertex names immaterial.
At one cut the only choice is which axiom to delete when both premises
have axiom producers. Either choice leaves one axiom between the two
companion occurrences, again with the same retained boundary incidence.
Correctness excludes a same-axiom loop. A nonaxiom principal cut has
just the specified corresponding-premise contraction.
These are all overlaps and establish local confluence.

Induct on the decreasing vertex count to obtain confluence.
For paths with first steps \(N\to U,N\to V\), local confluence gives
a common descendant \(W\). All of \(U,V\) have smaller measure.
Induction joins the first chosen endpoint to a descendant of \(W\)
inside \(U\), and the other to a descendant of \(W\) inside \(V\).
Induction below \(W\) joins these two descendants. A zero-step path
is immediate. Isomorphic representatives do not change the argument.
Thus all finite reductions join. Two cut-free normal forms can join
only by isomorphism, proving the stated uniqueness rather than
uniqueness of sequent proof trees.
\end{proof}
\end{annexitem}

